\parindent=0pt
\parskip=5pt
%%%%%%%%%%%%%%%%%%%%%%%%%%%%%%%%%%%%%%%%%%%%%%%
%%%%%%%% Information about using colors in LaTeX %%%%%%%%%%%%%%%%
%%%% https://www.overleaf.com/learn/latex/Using_colours_in_LaTeX %%%%%%%%
%%%%%%%%%%%%%%%%%%%%%%%%%%%%%%%%%%%%%%%%%%%%%%%
\usepackage[svgnames, dvipsnames, table]{xcolor}
%\usepackage[colorlinks = true, allcolors = Cerulean, pagebackref=true, colorlinks,]{hyperref}
%\renewcommand\backrefxxx[3]{%
 % \hyperlink{page.#1}{$\uparrow$#1}%
%}
%%%%%%%%%%%%%%%%%%%%%%%%%%%%%%%%%%%%%%%%%%%%%%%

\usepackage[left=3cm,right=3cm,top=1cm,bottom=3cm]{geometry} 
\usepackage[utf8]{inputenc}
\usepackage[spanish, es-tabla]{babel}
\usepackage{fancyhdr}
\usepackage{afterpage}
\usepackage[colorlinks = true, allcolors = Cerulean]{hyperref}
\usepackage{amsmath}
\usepackage{amsfonts}
\usepackage{amssymb}
\usepackage{graphicx}
\usepackage{sectsty}
\usepackage[table]{xcolor}
\usepackage{longtable}
\usepackage{float}
\usepackage{tablefootnote}
\usepackage{multicol}
\usepackage{multirow}
%\usepackage{biblatex}
%\addbibresource{Bibliografia-cursos.bib}
%\usepackage[sectionbib]{bibunits}
\usepackage[backend=biber, style=alphabetic]{biblatex}
%\usepackage[backend=biber,style=ieee,refsection=subsection]{biblatex} % Bibliografías por partes
\addbibresource{Biblio.bib}
%%%%%%%%%%%%%%%%%%%%%%%%%%%%%%%%%%%%%%%%%%%%%%%%%%
%%%%%%%%%% Definición de colores UCR %%%%%%%%%%%%%
%%%%%%%%%%%%%%%%%%%%%%%%%%%%%%%%%%%%%%%%%%%%%%%%%%
%% Estos colores se definen en las pp. 28-31 del %
%% manual de identidad visual de la UCR  %%%%%%%%%
%%%%%%%%%%%%%%%%%%%%%%%%%%%%%%%%%%%%%%%%%%%%%%%%%%
\definecolor{Celeste-UCR}{RGB}{0, 192, 243}
\definecolor{Azul-UCR}{RGB}{0, 93, 164}
\definecolor{Verde-UCR}{RGB}{109, 192, 103}
\definecolor{Amarillo-UCR}{RGB}{255, 224, 106}


%\chapterfont{\color{NavyBlue}}  % sets colour of chapters
\sectionfont{\color{Azul-UCR}}  % sets colour of sections
\subsectionfont{\fontsize{13.3}{13}\selectfont \color{Azul-UCR}}
\subsubsectionfont{\fontsize{13.3}{13}\selectfont \color{Azul-UCR}}


\newcommand{\N}{\mathbb{N}}
\newcommand{\Z}{\mathbb{Z}}
\newcommand{\C}{\mathbb{C}}
\newcommand{\Q}{\mathbb{Q}}
\newcommand{\I}{\mathbb{I}}
\newcommand{\R}{\mathbb{R}}
\newcommand{\op}[1]{\operatorname{#1}}
\newcommand{\bb}[1]{\mathbb{#1}}
\newcommand{\fr}[1]{\mathfrak{#1}}
\newcommand{\legendre}[2]{\left( \dfrac{#1}{#2} \right)}
\newcommand{\mcd}{\operatorname{mcd}}




%\title{Propuesta de programas de los cursos avanzados de análisis}
%\author{Comisión de Malla Curricular EMat}
%\date{\today}
